\documentclass[11pt]{article}
\usepackage[utf8]{inputenc}
\usepackage[T1]{fontenc}
\usepackage{geometry}
\geometry{a4paper, margin=1in}

\begin{document}

\begin{center}
\Large\bfseries Endorsement Request --- cs.AI Fallback
\end{center}

\vspace{0.3cm}

\noindent\textbf{From:} 20242100136@dcc.edu.cn \\
\textbf{To:} guo_jing@ynu.edu.cn（云南大学软件学院 郭竞老师） \\
\textbf{Subject:} Request for arXiv cs.AI endorsement --- AI 系统方向跨域手势识别可复现性研究

\vspace{0.3cm}
\hrule
\vspace{0.3cm}

\noindent 尊敬的郭老师：

您好！我是滇池学院（云南大学滇池学院）的访问研究人员。最近准备向 arXiv 投稿一篇关于 AI 系统（WiFi CSI 跨域手势识别）的论文，主投递类别在 cs.LG 与 cs.AI 之间权衡，故同时向两位老师咨询 endorsement 支持。

\vspace{0.2cm}

\noindent 了解到您在云南大学软件学院从事联邦学习、迁移学习、医学图像分析方向研究，与我们的跨域泛化议题有交集，故冒昧来信。

\vspace{0.2cm}

\noindent\textbf{论文核心：}

\begin{quote}
\small
本文是一篇关于 AI 系统（WiFi CSI 跨域手势识别）的可复现性研究。在严格的 18-fold × 5-seed paired 协议下，我们审计了 10 种近年提出的跨域方法，结果全部不显著（$|d|<0.3$, $p>0.28$）。我们也实证检验了 Liu 2025 (arXiv 2512.04521) 的 CBAM+ResNet18 系统——在 paired × multi-seed 协议下，标准 CBAM+ResNet18 跨域仅 28.35\%，远低于其报告的 97.61\%。我们的核心贡献是揭示了"AI 系统的报告提升常常被弱 paired × 单 seed 协议放大"。
\end{quote}

\vspace{0.2cm}

\noindent\textbf{与您研究的相关性：}
\begin{itemize}
\item 您做联邦学习 + 迁移学习，对"跨域泛化不鲁棒"这一现象应该有深度理解
\item 论文虽然主体在 cs.LG，但作为 AI 系统应用视角，cs.AI 也是合理的主投递类别
\item 如果 cs.LG endorsement 失败，我们会切换到 cs.AI primary
\end{itemize}

\vspace{0.2cm}

\noindent\textbf{endorsement 操作：}
\begin{enumerate}
\item 登录您的 arXiv 账号
\item 访问 https://arxiv.org/auth/endorsers
\item endorse others → 搜索我的 submission ID（投稿后补充）
\item 点击 endorse（cs.AI 类别）
\end{enumerate}

\vspace{0.2cm}

\noindent\textbf{补充材料：}
\begin{itemize}
\item \texttt{arxiv\_submission/paper\_arxiv.tex}（LaTeX 源文件）
\item \texttt{arxiv\_submission/paper\_arxiv.pdf}（编译后 PDF）
\end{itemize}

\vspace{0.2cm}

\noindent 如蒙支持，万分感谢。

\vspace{0.3cm}

\noindent 此致 \\
敬礼

\vspace{0.3cm}

\noindent [你的名字] \\
20242100136@dcc.edu.cn \\
滇池学院（云南大学滇池学院）

\end{document}
